% ==================================================
% CHAPTER 8: DEPLOYMENT AND OPERATIONS
% ==================================================

\chapter{Deployment and Operations}
\label{ch:deployment}

\section{Components}

The service is deployed at
\url{https://hostelswap.aakashchawla.space/login}, behind nginx, with
interactive API documentation at \texttt{/docs}. A deployment needs only
two things: one Python process and one PostgreSQL database.

\begin{itemize}
  \item \textbf{Application}: the \texttt{hostelswap} package run with
  Uvicorn. It serves the API and the compiled React interface from the
  same process, so no Node.js and no separate web server for static
  files are needed.
  \item \textbf{Database}: PostgreSQL, hosted on Supabase
  \cite{supabase}. Any PostgreSQL with the \texttt{gen\_random\_uuid()}
  function works; the expiry fix was verified on a local PostgreSQL~18.
\end{itemize}

\section{Configuration}

\begin{table}[H]
\centering
\small
\begin{tabularx}{\textwidth}{>{\ttfamily\raggedright\arraybackslash}p{0.25\textwidth} L}
\toprule
\textrm{\textbf{Setting}} & \textbf{Meaning} \\
\midrule
DATABASE\_URL & PostgreSQL connection string. Empty means no database: the demo cohort is served (development only). \\
APP\_ENV & \texttt{production} turns on \texttt{Secure} cookies and refuses to start without \texttt{DATABASE\_URL}. \\
\bottomrule
\end{tabularx}
\caption{Environment settings}
\label{tab:env}
\end{table}

Settings are read from the environment first and from a \texttt{.env}
file at the repository root second. An empty value counts as unset.
\texttt{.env}, private keys and certificates are ignored by Git.

\section{Setting Up}

\begin{lstlisting}[language=bash]
python -m venv .venv && source .venv/bin/activate
pip install -e "./code[dev]"
pytest -q                          # 148 tests; 4 more with TEST_DATABASE_URL
python db/migrate.py               # apply pending migrations
python db/seed_hostels.py --reset  # optional: 4 hostels, 1000 students
uvicorn hostelswap.api:app         # http://127.0.0.1:8000
\end{lstlisting}

The setup steps were checked on a fresh clone after they failed for a
new team member (week-6 journal); they now install from \texttt{./code},
where \texttt{pyproject.toml} lives.

\section{Releasing a Change}

\begin{enumerate}
  \item Work on a branch from \texttt{development}; CI runs the tests on
  Python 3.10--3.12 on every push and pull request.
  \item Merge into \texttt{development} by pull request, then into
  \texttt{main} for a release; pushing \texttt{main} rebuilds the
  documentation site.
  \item If the change touches the schema, add a new numbered file in
  \texttt{db/migrations/} and run \texttt{python db/migrate.py} against
  the production database \emph{before} the new code serves traffic.
  \texttt{--status} lists applied and pending migrations without
  changing anything.
  \item If the change touches the interface, rebuild it with
  \texttt{npm --prefix frontend run build} and commit the build, so the
  server needs no Node.js.
\end{enumerate}

The order in step 3 matters. The code from the expiry fix writes the
\texttt{settled} column added by migration \texttt{0002}, so running it
against a database without \texttt{0002} would make every reply to a
proposal fail.

\section{Operations}

\begin{itemize}
  \item \textbf{Health}: \texttt{/api/schema} reports the live database
  structure, which shows at once whether the database is reachable and
  migrated.
  \item \textbf{Running a round}: the warden opens a round, students
  register, and the warden presses ``lock and run''. The run's chains are
  stored, so later page loads only read them.
  \item \textbf{Expiry}: no scheduler is needed. Overdue proposals are
  expired when the next chain is offered and whenever someone replies
  to them, which are exactly the moments when an overdue proposal could
  otherwise block someone.
  \item \textbf{Test data}: \texttt{seed\_round.py} tops up an open round
  with an enrolled cohort and can be run again safely.
\end{itemize}
